\section{Marco teórico}

La simulación cuántica clásica de circuitos permite estudiar algoritmos, validar resultados y analizar costos computacionales antes de su ejecución en
hardware cuántico. Entre los modelos disponibles, la simulación basada en vectores de estado ocupa un lugar central por su generalidad, aunque dicha
generalidad viene acompañada de un crecimiento exponencial en memoria y tiempo. En este contexto, resulta necesario revisar los fundamentos de la
representación del estado cuántico, los principales enfoques de simulación clásica, los conceptos de compresión de datos científicos y los criterios
de exactitud numérica empleados para evaluar representaciones comprimidas y no comprimidas de un mismo estado.

\subsection{\normalsize{Representación del estado cuántico}}

En el modelo de circuitos cuánticos, el estado de un sistema de $n$ qubits se representa mediante un vector de dimensión $2^n$:
\[
    \ket{\psi}= \sum_{i=0}^{2^n-1} \alpha_i \ket{i},
\]
donde cada amplitud $\alpha_i \in \mathbb{C}$ satisface la condición de normalización
\[
    \sum_{i=0}^{2^n-1} |\alpha_i|^2 = 1.
\]

Esta formulación expresa al sistema como una combinación lineal de estados base del espacio de Hilbert. En términos computacionales, cada amplitud debe
almacenarse en memoria clásica cuando se utiliza una representación explícita del vector completo. Si se emplea doble precisión, cada amplitud
compleja requiere 16 bytes: 8 para la parte real y 8 para la parte imaginaria. El costo teórico de memoria queda entonces dado por:
\[
    M(n)= 2^n \times 16\ \text{bytes}.
\]

El crecimiento exponencial de $M(n)$ constituye una de las principales limitaciones de la simulación cuántica clásica. Para valores moderados de $n$,
el almacenamiento del vector de estado ya exige cantidades de memoria comparables con las de sistemas de cómputo de alto desempeño. Por esta razón, la
representación del estado no solo es un concepto matemático fundamental, sino también el origen del problema computacional tratado en simulación
cuántica a gran escala.

\subsection{\normalsize{Evolución unitaria y simulación tipo Schrödinger}}

La evolución de un sistema cuántico aislado se modela mediante operadores unitarios. Si $U$ representa una compuerta cuántica o una secuencia de
compuertas, el nuevo estado del sistema se obtiene como
\[
    \ket{\psi'} = U\ket{\psi}.
\]

La simulación tipo Schrödinger consiste en mantener explícitamente el vector de estado y aplicar sobre él las transformaciones unitarias del circuito.
Su principal ventaja es la generalidad: cualquier circuito que pueda expresarse en el modelo de compuertas puede, en principio, simularse sin imponer
restricciones sobre el grado de entrelazamiento o la forma del estado intermedio. Esta propiedad convierte a la simulación Schrödinger en una base
natural para contrastar representaciones alternativas del mismo estado cuántico.

Su principal limitación es que tanto el costo de memoria como una parte importante del costo temporal crecen exponencialmente con el número de qubits.
En consecuencia, cualquier estrategia orientada a reducir el uso de memoria debe analizarse sin perder de vista el formalismo de evolución unitaria que
define el problema original.

\subsection{\normalsize{Patrones de acceso inducidos por las compuertas cuánticas}}

Aunque el vector de estado contiene $2^n$ amplitudes, las operaciones elementales de un circuito no siempre requieren tratar todas las entradas de
manera arbitraria al mismo tiempo. Una compuerta de un solo qubit actúa sobre pares de amplitudes cuyos índices difieren en el bit asociado al qubit
objetivo. De manera similar, una compuerta de dos qubits actúa sobre grupos de cuatro amplitudes relacionadas por las combinaciones binarias de los
qubits involucrados.

Desde la perspectiva algorítmica, esto significa que la actualización del estado presenta patrones estructurados de acceso. La importancia de esta
observación reside en el hecho de que el costo práctico de simular un circuito depende
no solo del tamaño total del vector, sino también de cómo se distribuyen y reutilizan los accesos sobre sus componentes.

El concepto de localidad resulta entonces relevante. En computación, la localidad temporal describe la tendencia de ciertos datos a reutilizarse en un
intervalo corto de tiempo, mientras que la localidad espacial se refiere a accesos cercanos dentro de una misma región de memoria. Ambas nociones son
fundamentales para comprender por qué algunas representaciones del estado pueden explotarse con mayor eficiencia que otras.

En simulación tipo Schrödinger, la localidad no implica necesariamente que los elementos afectados se encuentren siempre contiguos en memoria, sino que
las compuertas inducen relaciones regulares entre índices. Dichas regularidades permiten razonar sobre patrones repetitivos de lectura y escritura, y
por ello resultan importantes cuando más adelante se analizan estrategias de almacenamiento para vectores de gran tamaño.

\subsection{\normalsize{Alternativas para reducir memoria en simulación cuántica clásica}}

La literatura ha propuesto diversos enfoques para reducir el costo espacial de la simulación cuántica. Las redes tensoriales y métodos relacionados
aprovechan estructuras favorables del circuito, como bajo entrelazamiento o geometrías particulares, para evitar la materialización explícita del
vector completo. De forma análoga, los diagramas de decisión cuánticos explotan redundancias algebraicas y simetrías que permiten representar ciertos
estados de manera compacta.

También existen técnicas basadas en partición del circuito, formulaciones tipo caminos de Feynman y esquemas híbridos espacio-tiempo, donde la memoria
se reduce al costo de incrementar el tiempo de cómputo o de restringir la clase de circuitos que pueden tratarse eficientemente. Estas alternativas
demuestran que el problema de memoria puede abordarse desde distintos formalismos, pero también evidencian que las ventajas obtenidas suelen depender
de propiedades estructurales específicas del circuito o del estado.

Cuando se desea conservar la representación explícita del vector de estado, otra línea de trabajo consiste en estudiar estrategias de almacenamiento
más eficientes sobre la misma formulación Schrödinger. En ese caso, el interés teórico deja de estar en sustituir el modelo de simulación y pasa a
concentrarse en cómo representar, transformar y recuperar los datos numéricos del estado sin alterar su significado físico.

Esta distinción es importante porque no todas las técnicas de reducción de memoria persiguen el mismo objetivo. Algunas cambian la representación
matemática del estado para aprovechar estructura; otras mantienen la misma representación y modifican únicamente la manera en que los datos se guardan,
se recuperan o se transfieren en memoria. Diferenciar ambos enfoques ayuda a interpretar con mayor claridad las decisiones presentadas en capítulos
posteriores.

\subsection{\normalsize{Compresión de datos científicos en punto flotante}}

La compresión de datos busca reducir el espacio requerido para almacenar o transmitir información. En términos generales, puede clasificarse en dos
grandes categorías: compresión sin pérdida y compresión con pérdida. La primera garantiza la reconstrucción exacta de los datos originales tras la
descompresión; la segunda acepta una perturbación controlada a cambio de mayores ratios de compresión.

En datos científicos representados en punto flotante, la distinción entre ambos enfoques es especialmente importante. Muchas aplicaciones numéricas
toleran errores acotados y, por ello, emplean compresión con pérdida. Sin embargo, cuando se requiere preservar exactamente los valores originales,
solo la compresión sin pérdida resulta admisible.

Los arreglos de punto flotante presentan desafíos particulares para la compresión porque con frecuencia exhiben alta entropía aparente. No obstante,
incluso en secuencias numéricas complejas pueden existir regularidades en su representación binaria, por ejemplo en signos, exponentes o patrones a
nivel de bytes. Por esta razón, una práctica común en compresión científica consiste en aplicar primero transformaciones reversibles que reorganicen
los datos y, posteriormente, utilizar compresores sin pérdida de propósito general o de alta velocidad.

Técnicas como \emph{shuffle} y \emph{bitshuffle} son ejemplos de este principio. Ambas reorganizan los bytes o los planos de bits de un conjunto de
valores para hacer más visibles ciertas regularidades al compresor posterior. Este tipo de preacondicionamiento es particularmente relevante en
arreglos homogéneos de gran tamaño, como los que aparecen en simulación numérica y cómputo científico.

Desde el punto de vista conceptual, estas transformaciones no cambian el significado de los datos, sino únicamente su organización antes de comprimir.
Su utilidad radica en que muchos compresores funcionan mejor cuando reciben secuencias con patrones repetitivos explícitos. En consecuencia, una parte
importante del problema no depende solo del compresor elegido, sino también de cuán favorable resulta la disposición binaria de los datos para dicho
compresor.

\subsection{\normalsize{Compresión por bloques y jerarquía de memoria}}

En el manejo de grandes volúmenes de datos científicos, una estrategia frecuente consiste en particionar los arreglos en bloques independientes. Este
principio aparece en técnicas de almacenamiento, procesamiento fuera de núcleo, cachés y bibliotecas de compresión orientadas a alto desempeño. La
motivación general es que operar sobre unidades más pequeñas permite equilibrar mejor el costo de acceso, el aprovechamiento de la localidad y el uso
de memoria temporal.

La granularidad del bloque introduce un compromiso clásico. Bloques grandes tienden a capturar más redundancia y pueden favorecer mejores ratios de
compresión, pero suelen incrementar la latencia de acceso cuando solo una fracción pequeña de los datos es necesaria. Bloques pequeños reducen el
costo de acceso localizado, aunque pueden degradar el rendimiento del compresor y aumentar el peso relativo del metadato y la administración.

Asociado a esta idea aparece el concepto de jerarquía de memoria. En muchos sistemas de cómputo, una parte de los datos permanece en almacenamiento
principal, mientras otra se mantiene temporalmente en niveles de acceso más rápido. Las decisiones sobre reemplazo, reutilización y permanencia de los
datos en estos niveles intermedios se relacionan con la noción general de caché. Dentro de este marco, políticas como \emph{Least Recently Used}
(LRU) constituyen mecanismos clásicos para explotar localidad temporal cuando el conjunto de datos activo excede la capacidad de memoria inmediata.

La intuición detrás de LRU es sencilla: si un dato fue utilizado recientemente, es razonable suponer que tiene mayor probabilidad de reutilizarse en el
corto plazo que otro que lleva más tiempo sin accederse. Por ello, cuando la capacidad disponible es limitada, la política expulsa primero el elemento
menos recientemente utilizado. Aunque esta regla no siempre produce la decisión óptima, su fundamento coincide con un principio ampliamente aceptado en
sistemas de memoria: muchos programas exhiben conjuntos de trabajo que permanecen activos durante intervalos acotados de ejecución.

En este sentido, el valor teórico de una política como LRU no depende de una implementación particular, sino de su relación con la noción de
localidad temporal. Comprender este vínculo facilita interpretar por qué la reutilización reciente de datos puede traducirse en menores costos de
acceso y, en términos generales, en una mejor utilización de memoria intermedia.

\subsection{\normalsize{Exactitud numérica y métricas de comparación}}

La evaluación de representaciones comprimidas y no comprimidas de un mismo estado requiere distinguir entre ahorro de memoria, costo temporal y
exactitud numérica. En cuanto a memoria, una métrica habitual es la razón de compresión:
\[
    CR = \frac{S_{\text{sin comp.}}}{S_{\text{comp.}}},
\]
donde $S_{\text{sin comp.}}$ representa el tamaño del estado en su forma no comprimida y $S_{\text{comp.}}$ el tamaño total bajo compresión.

En cuanto a desempeño, resulta natural medir el tiempo total de ejecución y la sobrecarga relativa frente a una simulación equivalente sin
compresión:
\[
    OH = \frac{T_{\text{comp}} - T_{\text{sin comp.}}}{T_{\text{sin comp.}}}.
\]

Cuando el objetivo es preservar exactamente las amplitudes originales, la corrección debe formularse en primer lugar como una condición de igualdad
exacta respecto a una simulación de referencia sin compresión. En términos ideales, esto significa verificar que
\[
    \alpha_i^{(\text{ref})} = \alpha_i^{(\text{comp})} \quad \forall i.
\]

Esta condición expresa de forma directa el significado de una estrategia sin pérdida: el estado recuperado debe coincidir exactamente con el estado de
referencia, no solo aproximarse a él. Cuando la simulación utiliza la misma precisión numérica y el mismo orden de operaciones, esta igualdad puede
incluso comprobarse a nivel bit a bit.

Como medida auxiliar de reporte, puede emplearse el error absoluto máximo:
\[
    e_{\text{abs}} = \max_i \left|\alpha_i^{(\text{ref})} - \alpha_i^{(\text{comp})}\right|.
\]
Esta métrica no sustituye el criterio de igualdad exacta, pero resulta útil para resumir numéricamente cualquier discrepancia en un único valor. En un
esquema estrictamente sin pérdida se espera que $e_{\text{abs}} = 0$.
